
\section{\BookRecordIndex{T0143}格罗莫夫--威滕不变量}\label{sec:1-4}
\zhindex{geluomofu-weitengbubianliang}{格罗莫夫--威滕不变量}

\BookRecordIndex{T0143}格罗莫夫--威滕不变量来自\BookRecordIndex{T0145}格罗莫夫--威滕伪循环
\[
  \operatorname{ev}:\mathcal M^*_{0,k}(A;J)\longrightarrow M^k
\]
与 $M^k$ 中余维互补的循环之交。更准确地，设
$a_1,\dots,a_k\in H^*(M)$ 是纯次数上同调类，并满足
\[
  m:=\sum_{j=1}^k\deg(a_j)=\mu(A,k).
\]
可以在 $M^k$ 中选择一个循环 $\alpha$，使它庞加莱对偶于上同调类
\[
  a:=\pi_1^*a_1\smile\dots\smile\pi_k^*a_k,
\]
并且按照第 6.5 节所精确定义的意义，$\alpha$ 与\BookRecordIndex{T0145}格罗莫夫--威滕伪循环
$\operatorname{ev}:\mathcal M^*_{0,k}(A;J)\to M^k$ \BookRecordIndex{T0281}强横截。于是把\BookRecordIndex{T0143}格罗莫夫--威滕
不变量定义为带定向的交数
\[
  \operatorname{GW}_{A,k}^M(a_1,\dots,a_k)=\operatorname{ev}\cdot\alpha.
\]
我们可以把 $\alpha$ 选成乘积 $\alpha_1\times\dots\times\alpha_k$，其中
$\alpha_i\subset M$ 庞加莱对偶于 $a_i$。在这种情形下，\BookRecordIndex{T0143}格罗莫夫--威滕不变量是同调类
$A$ 中分别与每个循环 $\alpha_1,\dots,\alpha_k$ 相交的 \BookRecordIndex{T0003}$J$-全纯曲线
$u:S^2\to M$ 的带定向计数。

若对 $i\in I$ \BookRecordIndex{T0150}固定标记点 $z_i$，相应不变量
$\operatorname{GW}_{A,k}^{M,I}$ 也用合适的\BookRecordIndex{T0283}求值映射与 $\alpha$ 的交来定义。作为第一步
近似，可以取该\BookRecordIndex{T0283}求值映射为 $\operatorname{ev}$ 在子空间
$\mathcal M^*_{0,k}(A;\mathbf w,J)$ 上的限制，其中对 $i\in I$ 令 $z_i:=w_i$。
注意，当 $k=3$ 时两种不变量相同：
\[
  \operatorname{GW}_{A,3}^M(a_1,a_2,a_3)
  =\operatorname{GW}_{A,3}^{M,\{1,2,3\}}(a_1,a_2,a_3).
\]
然而，当 $\#I>3$ 时，两者不同；见第 7.3 节。特别地，事实证明不变量
$\operatorname{GW}_{A,k}^{M,I}$ 具有非常有趣的形式性质，它们反映亏格零的
Grothendieck--Knudsen 空间之同调
$H_*\bigl(\overline{\mathcal M}_{0,k}\bigr)$ 上的 operad 结构。（下面还会谈到这一点。）

重要的一点是，这些不变量只依赖于 $(M,\omega)$ 的形变类，也就是 $\omega$ 在 $M$ 上
\BookRecordIndex{T0363}辛形式空间中所在的连通分支。该论断的证明使用\BookRecordIndex{T0370}虚模循环的存在性。本书只证明足够多的内容，
以说明当 $\omega$ 在\BookRecordIndex{T0058}半正形式中变化时不变量不变。维数不超过 $6$ 的流形上的每个\BookRecordIndex{T0363}辛形式
都是\BookRecordIndex{T0058}半正的，因此在这些维数中可得形变不变性。

第 7.1 节定义不变量 $\operatorname{GW}_{A,k}^M$ 并给出最初的一些例子。为了让希望从
那里开始阅读的读者能够进入，本节包含许多对前面结果的交叉引用，可在需要时回看技术细节。
带有若干\BookRecordIndex{T0150}固定标记点的不变量 $\operatorname{GW}_{A,k}^{M,I}$ 在第 7.3 节定义。第 7.4 节
以射影空间中\BookRecordIndex{T0386}有理曲线的计数为例说明这些思想，第 7.5 节则讨论所谓完整\BookRecordIndex{T0143}格罗莫夫--威滕不变量。

为解释这一点，注意当\BookRecordIndex{T0066}标记点数 $k\ge3$ 时有\BookRecordIndex{T0377}遗忘映射
\[
  \pi:\overline{\mathcal M}_{0,k}(A;J)
  \longrightarrow\overline{\mathcal M}_{0,k},
\]
这里 $\overline{\mathcal M}_{0,k}$ 是亏格为 $0$、带 $k$ 个\BookRecordIndex{T0066}标记点的\BookRecordIndex{T0339}稳定曲线组成的
Grothendieck--Knudsen \BookRecordIndex{T0246}模空间，映射 $\pi$ 把\BookRecordIndex{T0340}稳定映射 $[u,\mathbf z]$ 送到其底\BookRecordIndex{T0072}层\BookRecordIndex{T0339}稳定
曲线 $[\mathbf z]$。
\footnote{它也称为 Deligne--Mumford 空间。不过，后一组名字更恰当地属于高亏格定义域的
相应空间。}
考虑映射
\[
  \operatorname{ev}\times\pi:
  \overline{\mathcal M}_{0,k}(A;J)
  \longrightarrow M^k\times\overline{\mathcal M}_{0,k}.
\]
只有在很特殊的条件下它才是\BookRecordIndex{T0337}伪循环（见习题 6.7.13）。尽管如此，从这一观点考察
\BookRecordIndex{T0143}格罗莫夫--威滕不变量，可以非常有助于理解它们应当具有的性质。特别地，如第 6 章末尾
所述，对 $\overline{\mathcal M}_{0,k}$ 中每个同调类 $\beta$，应当能够定义一个不变量，
使它在类 $a\in H^*(M^k)$ 上的值
\[
  \operatorname{GW}_{A,k}^M(a;\beta)
\]
等于如下交数：取表示 $a$ 的庞加莱对偶的循环 $\alpha$，再将
$\operatorname{ev}$ 限制到 $\pi$ 下表示类 $\beta$ 的循环之原像上，所得映射与 $\alpha$
的交数。当 $\beta$ 是基本类 $[\overline{\mathcal M}_{0,k}]$ 时，重新得到不变量
$\operatorname{GW}_{A,k}^M$；当
$\beta=[\mathrm{pt}]\in H_0(\overline{\mathcal M}_{0,k};\mathbb Q)$ 时，得到所有\BookRecordIndex{T0066}标记点
固定的不变量 $\operatorname{GW}_{A,k}^{M,\{1,\dots,k\}}$。

这种观点的力量立即可见。如果用\BookRecordIndex{T0106}顶层 $\mathcal M_{0,k}$ 中一点 $\mathbf w$ 表示
$[\mathrm{pt}]\in H_0(\overline{\mathcal M}_{0,k})$，那么 $\mathbf w$ 是 $S^2$ 中 $k$ 个
两两不同的点组成的元组，不变量 $\operatorname{GW}_{A,k}^M(a;\mathrm{pt})$ 显然就是
$\operatorname{GW}_{A,k}^{M,\{1,\dots,k\}}$。然而，若取 $\mathbf w$ 位于另一个\BookRecordIndex{T0072}层中，
所计数的就是定义域以某棵固定树为模型的曲线。若不变量
$\operatorname{GW}_{A,k}^M(a;\beta)$ 确有良好定义，这两种计数必然相同。这就蕴含
\BookRecordIndex{T0143}格罗莫夫--威滕不变量应满足所谓的（分裂）公理。

第 7.5 节表述\BookRecordIndex{T0143}格罗莫夫--威滕不变量的 Kontsevich--Manin 公理，并用我们的不变量
$\operatorname{GW}_{A,k}^{M,I}$ 来解释它们。除（分裂）公理外，所有公理都容易建立。
由于证明该公理需要\BookRecordIndex{T0256}粘合定理，证明推迟到第 10.8 节。其中最重要的特殊情形是分解法则
\begin{align*}
  &\operatorname{GW}_{A,4}^{M,\{1,2,3,4\}}(a_1,a_2,a_3,a_4)\\
  &\quad=
  \sum_{B\in H_2(M)}\sum_{\nu,\mu}
  \operatorname{GW}_{A-B,3}^M(a_1,a_2,e_\nu)\,
  g^{\nu\mu}\,
  \operatorname{GW}_{B,3}^M(e_\mu,a_3,a_4),
\end{align*}
其中 $e_\nu$ 遍历有理上同调 $H^*(M)$ 的一组基，而 $g^{\nu\mu}$ 是配对矩阵
\[
  g_{\nu\mu}:=\int_M e_\nu\smile e_\mu
\]
的逆矩阵。这是证明\BookRecordIndex{T0228}量子杯积结合律的关键要素。

\paragraph{\BookRecordIndex{T0254}粘合。}
最后一个技术性很强的章节是第 10 章，它发展了 \BookRecordIndex{T0003}$J$-全纯曲线的\BookRecordIndex{T0254}粘合原理。可以把这一原理
看成格罗莫夫紧致性定理的逆命题。它断言：若两条或更多条 \BookRecordIndex{T0003}$J$-全纯曲线相交，并满足合适
的\BookRecordIndex{T0166}横截性条件，那么存在一列表示其同调类之和的 \BookRecordIndex{T0003}$J$-全纯曲线
$u_\alpha:S^2\to M$，在\BookRecordIndex{T0141}格罗莫夫收敛的意义下逼近它们。该定理最先由 Ruan--Tian
\srcite{345} 证明。我们给出另一种论证：它模仿规范理论中的构造，以仔细选择\BookRecordIndex{T0188}截断函数为
基础。本书的证明与 1994 年版本中的证明本质相同，但细节丰富得多。尽管建立\BookRecordIndex{T0257}粘合映射的
全部性质需要大量工作，底\BookRecordIndex{T0072}层思想并不太难理解；我们力求以便于理解的方式写出这一章。

本章所用技术在书中其他地方不再出现，因此把它推迟阅读不会造成严重后果。从第~\ref{ch:applications-symplectic-topology} 章可以
看出，\BookRecordIndex{T0003}$J$-全纯曲线有许多有趣应用并不使用\BookRecordIndex{T0254}粘合。这些应用以具有适当性质的 \BookRecordIndex{T0003}$J$-全纯曲线
之存在性为基础。不过，一旦需要计数曲线条数，几乎总要以（分裂）公理的形式使用\BookRecordIndex{T0254}粘合。
但往往又不需要该公理的全部力量。第 10.9 节证明了一个足以建立\BookRecordIndex{T0229}量子上同调的特殊情形。
